\answer{ex:21:1}{(a)~$T=\tau_{\text{rec}}\ln\frac{1}{1-x^*}$: $1.84$~s y $3.68$~s. (b)~De $3.36$ a $4.24$~s: la sensibilidad es $\dd T/\dd x^*=\tau_{\text{rec}}/(1-x^*)=80$~s; cuando $x^*\to1$, las ráfagas son irregulares.}
\answer{ex:21:2}{(a)~$8.6$~W, como en~\eqref{eq:energy}. (b)~$205$~Mbit/s, $0.2$~W: $2\cdot10^3$ veces más que el cultivo ($10^{-4}$~W).}
\answer{ex:21:3}{$x_{\text{est}}=1/(1+Ur_b\tau_{\text{rec}})<x^*$ cuando $r_b>(1/x^*-1)/(U\tau_{\text{rec}})=0.28$~Hz para $x^*=0.9$ y $0.025$~Hz para $x^*=0.99$; la fuerza de la sinapsis cae a $x_{\text{est}}$, es decir, un $10\%$ y un $1\%$.}
\answer{ex:21:4}{Las $u(t-k)$ son ortonormales, $m_k=\|P_Vu(t-k)\|^2$, donde $P_V$ es el proyector sobre el subespacio $N$-dimensional generado por las salidas; por la desigualdad de Bessel, $\sum_k m_k\le N$.}
\answer{ex:21:5}{$\mathrm{MC}\approx9$--$11$ con $\tanh$ y $15$--$18$ para el reservorio lineal (tres semillas), ambos $\le30$: la no linealidad se paga con memoria.}
\answer{ex:21:6}{(a)~$n\le102$. (b)~$5.6$.}
\answer{ex:21:7}{Problema abierto. Control clave: congelar la lectura y comprobar si la mejora se mantiene tras una pausa sin estimulación; si no, cambió el estado, no los pesos.}
